% concatenated from v1.tex
\documentclass[a4paper,12pt]{amsart}

\pdfoutput=1
\usepackage{setspace}
\setstretch{1.1}
\usepackage[left=22mm,head=30mm,bottom=30mm,foot=20mm,right=22mm]{geometry}
\usepackage{amsmath,amsthm,amsfonts,amssymb,mathtools,mathrsfs,url,bm,enumitem,stackengine}
\usepackage{newtxtext}
\usepackage[varvw]{newtxmath}
%\usepackage{showkeys}
%\usepackage[hiresbb]{graphicx}
%\usepackage{xcolor}
\allowdisplaybreaks
\usepackage{hyperref}
\hypersetup{
	colorlinks=true,
	linkcolor=blue,
	citecolor=red,%magenta,
	%filecolor=cyan,
	%urlcolor=cyan,
}
\usepackage{appendix}

\mathtoolsset{showonlyrefs=true}
\theoremstyle{definition}
\newtheorem{definition}{Definition}[section]
\newtheorem{rem}{Remark}[section]
\newtheorem{remark}{Remark}[section]
\newtheorem{question}{Question}[section]
\newtheorem{conjecture}{Conjecture}[section]
\newtheorem{problem}{Problem}[section]
\newtheorem{ass}{Assumption}[section]
\newtheorem{notation}{Notation}[section]
\theoremstyle{plain}
\newtheorem{thm}{Theorem}[section]
\newtheorem{theorem}{Theorem}[section]
\newtheorem{cor}{Corollary}[section]
\newtheorem{corollary}{Corollary}[section]
\newtheorem{prop}{Proposition}[section]
\newtheorem{proposition}{Proposition}[section]
\newtheorem{lem}{Lemma}[section]
\newtheorem{lemma}{Lemma}[section]
\def\R{\mathbb{R}}
\renewcommand{\Re}{\operatorname{Re}}
\renewcommand{\Im}{\operatorname{Im}}
\renewcommand{\phi}{\varphi}
\newcommand\restr[2]{{% we make the whole thing an ordinary symbol
		\left.\kern-\nulldelimiterspace % automatically resize the bar with \right
		#1 % the function
		\vphantom{\big|} % pretend it's a little taller at normal size
		\right|_{#2} % this is the delimiter
}}
\newcommand{\closure}{\overline}
\let\phi\varphi
\let\epsilon\varepsilon
\DeclareMathOperator{\sgn}{sgn}
\DeclareMathOperator{\dist}{dist}
\newcommand{\ts}{\textsuperscript}
\DeclareMathOperator{\Arctan}{Arctan}
\newcommand{\relmiddle}[1]{\mathrel{}\middle#1\mathrel{}}
\DeclareMathOperator{\spn}{span}

\usepackage{tikz-cd}
% from https://tex.stackexchange.com/a/216042/
\tikzset{
	symbol/.style={
		draw=none,
		every to/.append style={
			edge node={node [sloped, allow upside down, auto=false]{$#1$}}}
	}
}
\newcommand{\function}[5]{%
	\begin{tikzcd}[
		column sep=2em,
		row sep=1ex,
		ampersand replacement=\&
		]
		#1\colon \&[-3em]
		#2\vphantom{#3} \arrow[r] \&
		#3\vphantom{#2} \\
		\&
		#4\vphantom{#5} \arrow[u,symbol=\in] \arrow[r,mapsto] \&
		#5\vphantom{#4} \arrow[u,symbol=\in]
	\end{tikzcd}%
}

%\newcommand{\fscomment}[1]{\color{blue}}
% Sorry but I changed your command a little bit so that the first letter is appropriately printed.
% You can use it as \fscomment{Blah blah...}
\newcommand{\fscomment}[1]{{\color{blue}#1}}
\newcommand{\kaicomment}[1]{\textcolor{green}{#1}}
\newcommand{\GVHcomment}[1]{\textcolor{teal}{#1}}
%\usepackage[colorlinks=true, pdfstartview=FitV, linkcolor=blue, citecolor=blue, urlcolor=blue,pagebackref=false]{hyperref}
%%%%%%%% MACROS %%%%%%%%%%%%%%%%%%%%%%%%%%%%%%%%%%
\def\cantor{K_{\!\scriptscriptstyle\io}}
\def\p{\partial}
\def\no{\noindent}
\def\io{{\infty}}
\def\curl{\operatorname{curl}}
\def\diam{\operatorname{diam}}
\def\disc{\operatorname{disc}}
\def\range{\operatorname{ran}}
\def\diag{\operatorname{diag}}
\def\dist{\operatorname{dist}}
\def\re{\operatorname{Re}}
\def\card{\operatorname{Card}}
\def\im{\operatorname{Im}}
\def\Id{\operatorname{Id}}
\def\Lg{\operatorname{Log}}
\def\dive{\operatorname{div}}
\def\grad{\operatorname{grad}}
\def\moo{C^{\io}}
\def\mooc{C^{\io}_{\textit c}}
\def\moob{C^{\io}_{\textit b}}
\def\moc{C_{\text c}}
\def\N{\mathbb N}
\def\Z{\mathbb Z}
\def\Q{\mathbb Q}
\def\R{\mathbb R}
\def\C{\mathbb C}
\def\rot{\mathbf C}
\def\poscal#1#2{\langle#1,#2\rangle}
\def\proscald#1#2{<\!\!#1, #2\!\!>_{_{\hskip-4pt \R^{d}}}}
\def\proscal3#1#2{<\!\!#1, #2\!\!>_{_{{\hskip-4pt\R^{3}}}}}
\def\Poscal#1#2{\left\langle#1,#2\right\rangle}
\def\poscalr#1#2{\prec\!#1,\!#2\!\succ}
\def\inpro{\mathlarger{\mathlarger{\lrcorner}}}
\def\Poscalr#1#2{\mathlarger{\mathlarger{\prec}}#1,#2\mathlarger{\mathlarger{\succ}}}
\def\poi#1#2{\left\{#1,#2\right\}}
\def\NORM#1{\bigl\Vert#1\bigr\Vert}
\def\norm#1{\Vert#1\Vert}
\def\trinorm#1{\vert\hskip-1pt\Vert #1\Vert\hskip-1pt\vert}
\def\val#1{\vert#1\vert}
\def\Val#1{\left\vert#1\right\vert}
\def\valjp#1{\langle#1\rangle}
\def\diff#1#2{\frac{\partial #1}{\partial #2}}
\def\lap{\triangle}
\def\vu{ u }
\def\vn{ \nabla }

\def\l2{L^2(\R^{n})}
\def\L2{L^2(\R^{2n})}
\def\supp{\operatorname{Supp}}
\def\spec{\operatorname{spect}}
\def\singsupp{\operatorname{singsupp}}
\def\sign{\operatorname{sign}}
\def\tr#1{{^t}#1}
\def\tendinv#1#2#3#4{{{#1}\begin{matrix}\\ \longleftarrow\\ \hfill{\scriptstyle #3\rightarrow #4}\end{matrix}#2}}
\def\tend#1#2#3#4{{{#1}\begin{matrix}\\ \longrightarrow\\ \hfill{\scriptstyle #3\rightarrow #4}\end{matrix}#2}}
\def\rbar{\overline\R}
\def\rpbar{\overline\R_+}
\def\car#1{{\bold 1 }_{#1}}
\def\ps{pseudodifferential}
\def\op#1{{\text{Op}(#1)}}
\def\ops#1#2{{\text{Op}_{#1}(#2)}}
\def\w#1{{#1^{\text {Wick}}}}
\def\lloc#1{L^{#1}_{\text{loc}}}
\def\hloc#1{H^{#1}_{\text{loc}}}
\def\lcom#1{L^{#1}_{\text{comp}}}
\def\hcom#1{H^{#1}_{\text{comp}}}
\def\RZ{\R^{2n}}
\def\lip{\operatorname{Lip}}
\def\trace{\operatorname{trace}}
\def\sign{\operatorname{sign}}
\def\cps{condition $(\psi)$}
\def\hs{{\hskip15pt}}
\def\vs{\vskip.3cm}
\def\indic#1{\mathbf 1_{#1}}
\let \dis=\displaystyle
\let\no=\noindent
\let \dis=\displaystyle
\let\poo=\scriptstyle
\let\pooo=\scriptscriptstyle
\let\no=\noindent
\def\mat22#1#2#3#4{\begin{pmatrix}#1&#2\\ #3&#4\end{pmatrix}}
\def\matdu#1#2{\begin{pmatrix}#1\\ #2\end{pmatrix}}
\def\mattre#1#2#3{\begin{pmatrix}#1\\ #2\\#3\end{pmatrix}}
\def\rank{\operatorname{rank}}
\def\wt#1{\widetilde{#1}}
\def\ep{\varepsilon}

\def\mP{\mathcal{P}}
\def\mPl{\mathcal{P}^{\text{log}}} 
\def\Rt{\mathbb{R}^d}

\begin{document}
	
	
	\title[Liouville Theorems for SNS via the Radial Velocity Component]{Liouville Theorems for Stationary Navier-Stokes Equations via the Radial Velocity Component}

%\title[Liouville-Type Results for SNS under radial integrability assumptions]{Liouville theorems for stationary Navier-Stokes Equations under radial integrability assumptions}
	
	\author{Gast\'{o}n Vergara-Hermosilla}
\address{Institute for Theoretical Sciences, Westlake University, 310030 Hangzhou, Zhejiang, People's
Republic of China}
\email{gaston.v-h@outlook.com}	
	
	\maketitle
	
	
\begin{abstract}
We study Liouville-type results for the stationary Navier--Stokes equations in $\mathbb{R}^3$. We prove that any $\dot{H}^1(\mathbb{R}^3)$ solution is trivial under an integrability condition imposed only on the radial component of the velocity, namely $u_\rho(x) \in L^p(\mathbb{R}^3)$ with $3/2 < p \leq 3$. We also establish a uniqueness result in a variable-exponent setting, where an $L^6$-type condition is required only on a bounded region, while the exponent approaches the critical value $3$ at infinity. Our analysis reveals that the rigidity of the stationary Navier--Stokes system can be driven by localized and  radial  integrability properties, rather than uniform global conditions.
\end{abstract}

%---------------------------------------------------
 \section{Introduction and main results}
 
 	In this note we consider the stationary Navier-Stokes equations 
	\begin{equation}\label{SNS}
		\begin{cases}
			-\Delta u+  ( u\cdot \nabla ) u +\nabla \pi  =0,
			\\
			\quad \nabla \cdot u=0, 
		\end{cases}
		\quad 
		\text{in }\R^3,
	\end{equation} 
 where $\vu $ stands the velocity vector field and $\pi  $ denotes the  pressure. The construction of solutions $(\vu, \pi)$ for this equation in the spaces $\dot{H}^1(\mathbb{R}^3) \times \dot{H}^{\frac{1}{2}}(\R^3) $  is classical (see \cite[Theorem 16.2]{lemarie2016navier}), however  uniqueness of solutions is  until now a challenging and open problem. This motivates the following conjecture, which has been originally proposed by G. Galdi in \cite[Remark X.9.4 and Theorem X.9.5]{galdi2011introduction} and by G. Seregin in \cite{Ser2016}):
\begin{conjecture}
Show that any solution $\vu$ of (\ref{SNS}) fulfilling the conditions
\begin{equation}\label{Conjecture}
\vu \in \dot{H}^1(\R^3) \qquad\mbox{and}\qquad \vu(x)\to 0 \mbox{ as } |x|\to +\infty,
\end{equation}
is identically equal to zero.
\end{conjecture}
Note that the Liouville problem in the two-dimensional case was established by Gilbarg and Weinberger in \cite{GW78}, while the four-dimensional case was proved by G. Galdi in \cite{galdi2011introduction}. The three-dimensional case, however, remains widely open and has attracted significant interest within the mathematical fluid mechanics community.
On the other hand, by Sobolev embedding, any solution
 $u \in \dot{H}^1(\mathbb{R}^3)$ also belongs to $L^6(\mathbb{R}^3)$ which already implies a certain decay as $x\to \infty$. However, this information alone does not appear to be sufficient to deduce the triviality of the solution. Over the years, several partial results related to Conjecture \ref{Conjecture} have been obtained, showing that additional integrability or structural assumptions ensure that $u=0$. In this direction  G.~Galdi  proved in \cite{galdi2011introduction} that $u \in L^{9/2}(\mathbb{R}^3)$ in fact implies  that $u \equiv 0$. This result was later refined by D.~Chae and J.~Wolf in \cite{ChaeWolf}, who obtained a logarithmic improvement of the $9/2$ condition.    N.~Lerner observed in \cite{Lerner26} that the global $L^{9/2}$ assumption can be relaxed by separating low and high frequencies. More precisely, he showed that it suffices to assume $u_{[0]} \in L^{9/2}(\R^3)$, with  $u_{[0]}$ denoting a projection of $u$ onto the subspace of vector fields whose Fourier support contains a neighborhood of the origin. More recently, the author of this paper  proved in \cite{V2026} that triviality already follows under assumptions of the form $u \in L^{9/2 + \varepsilon(\cdot)}(\mathbb{R}^3)$, where $\varepsilon(\cdot)>0$.  For more results about Liouville type theorems, we refer to the reader to Chapter 16 in the book of P.G. Lemari\'e-Rieusset \cite{lemarie2016navier} and the references therein.\\ 

To the best of the author’s knowledge, most Liouville-type results in the literature are isotropic, i.e., invariant under rotations of the coordinate system, while anisotropic results remain scarce. A seminal work in this  direction is due to D. Chae \cite{Chae23}, who proved uniqueness by considering sufficient conditions depending on the direction of coordinates and each component of the velocity on mixed norm Lebesgue spaces. This results was extended by Z. Zhang and Q. Zu in \cite{ZZ25}. 
In this paper, we study new sufficient conditions for Liouville-type theorems that depend explicitly only on  the radial component of the velocity field, thereby offering a different perspective on this problem.
 
We have to new theorems:
  %%%%%%%%%%%%%%%%%%%%%%%%%%%%%%
 \begin{theorem}\label{thm1} 
 Let $\vu  \in \dot{H}^1(\mathbb{R}^3) $ be a solution of (\ref{SNS})
 and let $u_\rho(x)$ be the radial component of such solution in spherical coordinates. 
 Suppose in addition that
 $u_\rho \in L^{p}(\mathbb{R}^3)$ with $3/2<p\leq 3$ .
 Then, $u $ is identically equal to zero.
\end{theorem}

\begin{remark}
This result establishes the uniqueness of steady solutions under a weaker additional integrability assumption compared to known isotropic and anisotropic Liouville-type results for the stationary Navier--Stokes equations. 
More precisely, Theorem \ref{thm1} imposes sufficient conditions only on the radial component of the velocity field, without requiring any restrictions on the remaining two components.
\end{remark}

\begin{remark}
An interesting case in the Liouville theorem of D. Chae in \cite{Chae23} corresponds to the assumption (where 
$\tilde x_1 = (x_2, x_3), \ \tilde x_2= (x_1, x_3), \ \tilde x_3 = (x_1, x_2)$) 
\[
u_i \in \big(L_{x_i}^{\frac{3}{2}} \cap L_{x_i}^{6}\big) L_{\tilde x_i}^{6}(\mathbb{R} \times \mathbb{R}^2), 
\quad \text{for each cartesian coordinate indexed with } i = 1,2,3.
\]
This condition reflects an anisotropic decay, combining  mild integrability 
in planar directions with stronger decay in the orthogonal direction, and is sufficient 
to ensure the triviality of the solution.
By contrast, Theorem \ref{thm1} relies on the assumption 
$u_\rho \in L^{p}(\mathbb{R}^3)$ with $3/2 < p \leq 3$. In particular, this involves 
a range of integrability exponents, and only on a single component 
of the velocity field. 
%This suggests that, within this framework, the nonlinear structure can be controlled under  weaker integrability assumptions.
\end{remark}

\begin{remark}
A particular case of our result corresponds to the assumption 
$u_\rho \in L^{3}(\mathbb{R}^3)$. At a heuristic level, this can be compared with the classical 
a priori condition $u \in L^{9/2}(\mathbb{R}^3)$ due to G.~Galdi 
\cite{galdi2011introduction};   
the latter integrability exponent naturally arises by applying Hölder's inequality 
to an integral term involving a cubic power of $|u|$, 
which leads to the balance condition
\[
\frac{2}{3} = \frac{3}{p},
\]
whose unique solution is $p = \frac{9}{2}$.
In contrast, our approach relies on a different distribution of the analogous integral term, 
involving two factors of the full velocity field $u$ and one factor of its radial 
component $u_\rho$. This leads, via Hölder's inequality, to the relation
\[
\frac{2}{3} = \frac{2}{p_1} + \frac{1}{p_2}.
\]
A natural and convenient choice of exponents is given by $p_1 = 6$ and $p_2 = 3$. 
This selection is consistent with the Sobolev embedding 
$\dot H^1(\mathbb{R}^3) \hookrightarrow L^6(\mathbb{R}^3)$, 
which provides control of the velocity field in $L^6$, while the radial 
component is handled in $L^3$. 
This argument also plays a role in the proof of our theorem, specifically in the control of the decay of terms involving the head pressure 
$Q = \pi + \frac{|u|^2}{2}$ (see \eqref{forheuristic1}--\eqref{forheuristic2} for further details).
\end{remark}
 


 


In our second main result we establish a  Liouville-type result that improves upon $L^{3}$ condition for the radial component of the velocity field. More precisely, we show that the triviality of solutions already follows from a weaker integrability assumption of the form $u \in L^{3+\eta}(\mathbb{R}^3)$, with $\eta>0$.  This theorem reads as follows.

% \newpage 
  \begin{theorem}\label{thm1.2} 
 Let $\vu  \in \dot{H}^1(\mathbb{R}^3) $ be a solution of (\ref{SNS})
 and let $u_\rho(x)$ be the radial component of such solution in spherical coordinates. 
Consider $R_0>3$ fixed  and let $\eta(\cdot)$ be a scalar function defined by  $\eta(x)=  3 $ for $|x|< R_0$ and $\eta(x)=   \frac{3R_0}{|x|},$ for $|x| \geq  R_0$.  
 Suppose in addition that
 $u_\rho \in L^{3 + \eta(\cdot) }(\mathbb{R}^3)$.
 Then, $u $ is identically equal to zero.
\end{theorem} 

\noindent We now comment on this result and its relation to the existing literature.

\begin{remark} 
Theorem \ref{thm1.2} offers an alternative perspective toward uniqueness: 
instead of imposing a global $L^6$ condition\footnote{As mentioned above, it remains an open question whether the global condition 
$u \in L^6(\mathbb{R}^3)$ implies the triviality of solutions.} on $\mathbb{R}^3$, 
we require such integrability only on a fixed compact domain 
(note that $\eta(x)=3$ on $B(0,R_0)$), while allowing the exponent to 
decrease radially toward $3$ as $|x| \to \infty$. 
This formulation captures a transition from the desired  integrability regime 
in a bounded region to the $L^3$ behavior attained asymptotically at infinity.
\end{remark}

\begin{remark}
Theorem \ref{thm1.2} also shows that the integrability condition can be effectively 
localized at infinity. More precisely, the critical exponent $3$ is only required 
asymptotically, and may be approached from above at a controlled rate. 
We further emphasize that the radius $R_0$ is arbitrary, so the region where stronger 
integrability is imposed can be taken arbitrarily far from the origin.
\end{remark}


To prove Theorem \ref{thm1.2}, we rely on a more general uniqueness result in the context of Lebesgue spaces with variable exponents. This functional setup allows us to capture different integrability behaviors within a single space. Such spaces have already been used to prove Liouville-type theorems for the stationary Navier–Stokes equations and to overcome the well-known $9/2$ exponent threshold, we refer the interested reader to \cite{CV2023} for more details.\\

The rest of the paper is structured as follows. Section \ref{sect:2} gives a brief overview of variable exponent Lebesgue spaces: their definitions, main properties, and a few useful lemmas. Section \ref{sec:3} then presents the proof of our main results.


%%%%%%%%%%%%%
\section{Preliminaries}\label{sect:2}

We begin this section by introducing some notations used throughout this paper; given $R>1$, we consider 
\[
C( R) := \left\{ x \in \mathbb{R}^3 \ \big | \ \frac R2 < |x| < R \right\}.
\]

For clarity of the presentation, in this section we recall 
several definitions and properties of variable Lebesgue spaces. Given $\mathcal{X}\subseteq \mathbb{R}^{n}$, let $\mathcal{P}(\mathcal{X})$ be the set of measurable functions $p(\cdot):\mathcal{X}\rightarrow[1,+\infty]$.  The elements of $\mathcal{P}(\mathcal{X})$ are called variable exponents. For $p(\cdot)\in \mathcal{P}(\mathcal{X})$, we consider the following notations
$$
 p^{-}:=\operatorname{essinf}_{x\in\mathcal{X}}p(x), \ \  p^{+}:=\operatorname{esssup}_{x\in\mathcal{X}}p(x).
$$
In what follows we will always  consider  $1<p^-\leq p^+<+\infty$.\\

Given  $\mathcal{X}\subseteq\mathbb{R}^{n}$,  $p(\cdot)\in \mathcal{P}(\mathcal{X})$,  and a measurable function $f(x)$,  we set the quantity
\begin{equation}\label{equation2.2}
   \|f\|_{L^{p(\cdot)}}(\mathcal{X}):=\inf\left\{\lambda>0: 
   \int_\mathcal{X} \left|\frac{f(x)}{\lambda}\right|^{p(x)}dx
\leq 1\right\}.
\end{equation}
If the set involved in  \eqref{equation2.2} is empty,  then we define $\|f\|_{L^{p(\cdot)}} = \infty$.
Note that,  if the exponent function $p(\cdot)$ is a constant, i.e. if $p(\cdot)=p\in [1,\infty)$, then we obtain the usual norm of Lebesgue spaces.
\smallbreak

\begin{definition}\label{de2.2}
Given  $\mathcal{X}\subseteq\mathbb{R}^{n}$ and $p(\cdot)\in \mathcal{P}(\mathcal{X})$,  the variable exponent Lebesgue space $L^{p(\cdot)}(\mathcal{X})$ is defined as the set of  measurable functions $f(x)$ satisfying  $\|f\|_{L^{p(\cdot)}(\mathcal{X})}<\infty$.
\end{definition}
At this point we must stress the fact that the speaces $L^{p(\cdot)}(\mathcal{X})$ are in fact Banach spaces associated with the norm $\|\cdot\|_{L^{p(\cdot)}}$.
In the following, we state some of their main properties.%--------------------------------
\begin{lemma}[H\"{o}lder inequality]
Consider a domain $\mathcal{X}\subseteq \mathbb{R}^n$ and $q(\cdot),\,r(\cdot),\,p(\cdot)\in \mathcal{P}(\mathbb{R}^n)$ such that 
$\frac{1}{p(x)}=\frac{1}{q(x)}+\frac{1}{r(x)}$, for $x\in \mathcal{X}$. 
Then, given $u\in L^{q(\cdot)}(\mathcal{X}) $ and $v \in L^{r(\cdot)}(\mathcal{X})$, the pointwise product $uv$ belongs to  $L^{p(\cdot)}(\mathcal{X})$, and there exists  a positive constant $C$ such that 
\begin{equation}\label{eq2.3}
\|uv\|_{L^{p(\cdot)}(\mathcal{X})} \leq C\|u\|_{L^{q(\cdot)}(\mathcal{X})}\|v\|_{L^{r(\cdot)}(\mathcal{X})}.
\end{equation}
\end{lemma}
%--------------------------------
A proof of this lemma can be consulted in  \cite[Lemma 3.2.20]{Diening_Libro}.
%%%%%%%%%%%%%%%%%%%%%%%%%%%%%%%%%%%%%%%%%%%%%%%%%%%
\begin{definition}\label{globallogholder}
Consider  a domain $\mathcal{X}\subseteq \mathbb{R}^d$ 
and  $p(\cdot)\in \mathcal{P}(\mathcal{X})$.
The variable exponent $ p(\cdot): \mathcal{X} \to \mathbb{R} $ is  called locally log-Hölder continuous on $ \mathcal{X} $ if there exists $ C_1 > 0 $ such that
\begin{equation}\label{relation_loc_logholder}
|p(x) - p(y)| \leq \frac{C_1}{\log(e + 1/|x - y|)},
\end{equation}
for all $ x, y \in \mathcal{X} $. We say that $p(\cdot) $ satisfies the  log-Hölder decay condition if there exist a constants $ p_\infty$ and  $ C_2 > 0 $ such that, for all $ x \in \mathcal{X} $  
\begin{equation}\label{relation_decay_logholder}
|p(x) -  p_\infty| \leq \frac{C_2}{\log(e + |x|)}
.
\end{equation}
The variable exponent $ p(\cdot) $ is  called globally log-Hölder continuous in $ \mathcal{X} $ if it is locally log-Hölder continuous and satisfies the log-Hölder decay condition, and this  class  is denoted by 
$\mathcal{P}^{log}(\mathcal{X})$. 
\end{definition}
%%%%%%%%%%%%%%%%%%%%%%%%%%%%%%%%%%%%%%%%%%%%%%%%%%%
Given a  $ \mathcal{X} \subset \Rt$  measurable and   $p(\cdot) \in \mathcal{P}(\Rt)$,  the notation 
$p_\mathcal{X}(\cdot)$ stands for the variable exponent restricted to the set $\mathcal{X}$, \emph{i.e.} ${p}_\mathcal{X}(\cdot)=  p  (\cdot)_{|_{\mathcal{X}}}$. 
%%%%%%%%%%%%%%%%%%%%%%%%%%%%%%%%%%%%%%%%%%%%%%%%%%%
\begin{lemma}\label{PropositionLpplusminus}
Consider a measurable set $\mathcal{X}\subset \R^3$ and $p(\cdot)\in \mathcal{P}(\Rt)$ a variable exponent, assume that we have $|\mathcal{X}|$ has finite measure.  Then 
$$\|1\|_{L^{p_{\mathcal{X}}(\cdot) }(\mathcal{X})}\leq 2\max\{|\mathcal{X}|^{\frac 1{ p^-}},|\mathcal{X}|^{\frac 1{p^+}}\}.$$
\end{lemma}
%%%%%%%%%%%%%%%%%%%%%%%%%%%%%%%%%%%%%%%%%%%%%%%%%%%
A proof of this lemma can be found in \cite[Lemma 3.2.12]{Diening_Libro}.  
%%%%%%%%%%%%%%%%%%%%%%%%%%%%%%%%%%%%%%%%%%%%%%%%%%%
\begin{lemma}\label{lemmaLinftyLpvariable}
Let $\mathcal{X}\subseteq \R^3$ and $p(\cdot)\in \mathcal{P}(\mathbb{R}^3)$ a variable exponent. Then, we have the space inclusion $L^{\infty} (\mathcal{X}) \subset L^{p_{\mathcal{X}}(\cdot)} (\mathcal{X})$, if and only if $1\in  L^{p_{\mathcal{X}}(\cdot)} (\mathcal{X}) $ and  the following estimate follows
$$\|f\|_{L^{p_{\mathcal{X}}(\cdot)} (\mathcal{X})}\leq \|f\|_{L^\infty(\mathcal{X})}\|1\|_{L^{p_{\mathcal{X}}(\cdot)} (\mathcal{X})}.$$
In particular, the embedding holds if $|\mathcal{X}|$ has finite measure. 
\end{lemma}
%%%%%%%%%%%%%%%%%%%%%%%%%%%%%%%%%%%%%%%%%%%%%%%%%%%
The proof of this lemma can be consulted in the book \cite[Proposition 2.43]{CRUZ}. 
The next proposition presents relations between the norm  and modular functions associated to functions in $L^{p(\cdot)}$. 
%
\begin{proposition} \label{proposition.utendzero}
Consider a domain $\mathcal{X} \subseteq \R^n$ and  a variable exponent $p(\cdot) \in \mathcal{P}(\mathcal{X})$. If $\|f\|_{L^{p(\cdot)}(\mathcal{X})} > 1$, then we have 
\[
 \|f\|_{L^{p(\cdot)}(\mathcal{X})} \leq \left(\int_{\mathcal{X}} |f(x)|^{p(x)}dx \right)^{1/p_-}.
\]
On the other hand, if we have $\|f\|_{L^{p(\cdot)}(\mathcal{X})} \leq 1$, then
\[
 \|f\|_{L^{p(\cdot)}(\mathcal{X})} \leq \left(\int_{\mathcal{X}} |f(x)|^{p(x)}dx \right)^{1/p_+}.
\] 
\end{proposition}
%
\noindent The proof of this proposition can be found in \cite[Chapter 2, page 25]{CRUZ}.  
The following result will be important for proving our main results. Although the proof is standard, we will present an additional proof for clarity of the presentation.
 
 %
\begin{lemma} \label{Lemma.utendzero}
Let $p(\cdot)\in \mathcal{P}(\R^3 )$  and $f\in L^{p(\cdot)}(\R^3 )$. 
Then, 
\begin{equation}
\lim_{R\to +\infty} \| f\|_{L^{p(\cdot)}(\mathcal{C}( R)   )} =0
.
\end{equation}
\end{lemma}
%
\begin{proof}
To begin note that, given $R>1$ and $\mathcal{C}( R)   \subset \mathcal{X}_R := \{ x\in \R^3 : \frac R 2 < |x| \}$, we can write 
\begin{equation}
\int_{\mathcal{C}( R)  } |f(x)|^{p(x)}dx
\leq 
\int_{\mathcal{X}_R} |f(x)|^{p(x)}dx<\infty.
\end{equation}
The, considering that  $\mathcal{X}_R \downarrow \emptyset \text{ as } R \to \infty $ and $|f(x)|^{p(x)}$ belongs to $L^1(\mathbb{R}^3)$, we conclude (using dominated convergence theorem)
$ \int_{\mathcal{X}_R} |f(x)|^{p(x)}dx \to  0\   
 \text{ as } R \to \infty .
$
Hence, we can write 
$$
\lim_{R \to \infty}
\int_{\mathcal{C}( R)  } |f(x)|^{p(x)}dx 
=0.$$
 
%Let $f(x):=|u(x)|^{p(x)} \in L^1(\mathbb{R}^3)$ and define 
%$f_R(x):=f(x)\chi_{E_R}(x)$. Then $f_R(x)\to 0$ pointwise and 
%$|f_R(x)|\le f(x)\in L^1(\mathbb{R}^3)$. Hence, by the dominated 
%convergence theorem,
%\[
%\lim_{R\to\infty}\int_{E_R}|u(x)|^{p(x)}dx=\lim_{R\to\infty}\int_{\mathbb{R}^3} f_R(x)\,dx=0.\]
 

%Thus, considering this fact and Proposition \ref{proposition.utendzero}, we conclude the desired limit.  
 % Thus, since
%\[\int_{\mathcal{C}( R)  } |u(x)|^{p(x)}dx \to 0,\]
\noindent 
%Thus, there exists $R_0>0$ such that for all $R\ge R_0$,
Now, note that, there exists $R_0>0$ such that for each $R \geq R_0$,
$
\int_{\mathcal{C}( R)  } |f(x)|^{p(x)}dx <1.
$
Thus, considering Proposition \ref{proposition.utendzero}, we get 
\[
\|u\|_{L^{p(\cdot)}(\mathcal{C}( R)  )}
\le \left(\int_{\mathcal{C}( R)  } |f(x)|^{p(x)}dx \right)^{1/p_+}
\quad \forall  R \ge R_0.
\]
Then,  by passing to the limit as $R\to\infty$, we conclude the desired limit.
 \end{proof} 

To continue, we present  a key lemmas that will be used in the proof of Theorem \ref{thm1.2}.

\begin{lemma}\label{prop:main_limit_thmover_new}
Consider $p(\cdot)$ be the variable exponent defined in Theorem \ref{thm1.2}, $R_0>3$ and $R\geq 2R_0^2$. Let  
$h:[1,\infty) \to [0,\infty)$ be such that 
$
h(R) \to 0  \text{ as } R \to \infty.
$ 
Then
\[
\lim_{R\to\infty} R^{1 - \frac{3}{p_{\mathcal{C}_R}^+}} h(R) = 0.
\]
\end{lemma}
\begin{proof} 
Let $x\in\mathcal{C}( R) $. 
Since $p(\cdot)$ is continuous and radially decreasing, there exists a decreasing and continuous function
$\tilde{p} : [0, \infty) \to [3, 6]$ such that $ p(x) = \tilde{p}(|x|) $ for all $ x \in \mathbb{R}^3 $.
For each $ x \in\mathcal{C}( R)  $, we have $|x| > R/2$, hence
\[
p(x) = \tilde{p}(|x|) \leq \tilde{p}(R/2),
\]
and therefore we can  write
\[
\operatorname{ess\,sup}_{x \in\mathcal{C}( R) } p(x) \leq \tilde{p}(R/2).
\]
Consider $\varepsilon > 0$. By continuity of $\tilde{p}$ at $R/2$, there exists $\delta > 0$ such that
\[
r \in (R/2, R/2 + \delta) \;\Rightarrow\; |\tilde{p}(r) - \tilde{p}(R/2)| < \varepsilon,
\]
which implies $\tilde{p}(r) > \tilde{p}(R/2) - \varepsilon$.
To continue, we define
\[
A(\delta,R) := \{x \in \mathbb{R}^3 : R/2 < |x| < R/2 + \delta\}.
\]
Then $A(\delta,R) \subset\mathcal{C}( R) $, and for all $x \in A_\delta$, we have
\[
p(x) > \tilde{p}(R/2) - \varepsilon.
\]
Moreover, since set $A(\delta,R)$ has positive measure,   we can write
\[
\operatorname{ess\,sup}_{x \in\mathcal{C}( R) } p(x) \geq \tilde{p}(R/2) - \varepsilon.
\]
Since $\varepsilon>0$ is arbitrary, we conclude
$
p_{\mathcal{C}_R}^+ = \tilde{p}(R/2).
$ 
Now since $R \geq 2R_0^2$, we have $R/2 \geq R_0^2$, and by assumption, we have 
\[
\left| \tilde p(R/2) - 3 \right| \leq \frac{C}{R/2} = \frac{2C}{R}<1.
\]
Defining $\varepsilon_R := \tilde p(R/2) - 3$, we conclude  
that there exists a sequence $(\varepsilon_R)$ such that
\[
p_{\mathcal{C}_R}^+ = 3 + \varepsilon_R, 
\quad 0<\varepsilon_R \le \frac{2C}{R}<1.
\]
In particular, note that 
\[
p_{\mathcal{C}_R}^+ = 3 + O\!\left( \frac{1}{R}\right).
\]

To continue, define $
\omega_R := 1 - \frac{3}{p_{\mathcal{C}_R}^+}$.   Thus, we can write 
\[
\frac{3}{p_{\mathcal{C}_R}^+}
= \frac{3}{3 + \varepsilon_R}
= \frac{1}{1 + \frac{\varepsilon_R}{3}}.
\]
Now, by using the Taylor expansion for $t = \varepsilon_R/3 < 1$, we get 
\[
\frac{1}{1+t} = 1 - t + O(t^2),
\]
and then we conclude 
\[
\frac{3}{p_{\mathcal{C}_R}^+}
= 1 - \frac{\varepsilon_R}{3} + O(\varepsilon_R^2).
\]
Therefore, we can write 
\[
\omega_R
= 1 - \frac{3}{p_{\mathcal{C}_R}^+}
= \frac{\varepsilon_R}{3} + O(\varepsilon_R^2)
= O\!\left(\frac{1}{R}\right).
\]
Now, using the identity 
$
R^{\omega_R} = \exp(\omega_R \ln R),
$ 
and the fact that $\omega_R = O(1/R)$, we get
\[
\omega_R \ln R = O\!\left(\frac{\ln R}{R}\right) \to 0.
\]
Thus, considering $s:=\omega_R \ln R$ and using the expansion $e^s = 1 + O(s)$ as $s \to 0$, we conclude
\[
R^{\omega_R} = 1 + O\!\left(\frac{\ln R}{R}\right).
\]
Provided with this information   we can write
\[
R^{1 - \frac{3}{p_{\mathcal{C}_R}^+}} = R^{\omega_R}
= 1 + \alpha_R,
\]
where $\alpha_R = O\!\left(\frac{\ln R}{R}\right)$.
Thus, we obtain 
\[
R^{1 - \frac{3}{p_{\mathcal{C}_R}^+}} h(R)
= (1 + \alpha_R) h(R)
= h(R) + \alpha_R h(R).
\]
Since $h(R) \to 0$ and $(\alpha_R)$ is bounded for $R \geq 1$, we have
\[
|\alpha_R h(R)| \le C |h(R)| \to 0 \quad \text{as } R \to \infty.
\]
Thus, conclude 
$ 
R^{1 - \frac{3}{p_{\mathcal{C}_R}^+}} g(R) \to 0   \text{ as } R \to \infty,
$ 
and we finish the proof. 
\end{proof}
%%%%%%%%%%%%%%%%%%%
\section{Proofs of main theorems}\label{sec:3} 
\begin{proof}[Proof of Theorem \ref{thm1}]
To begin note that, since $\vu  \in \dot{H}^1(\mathbb{R}^3) $, we have $\vu \in L^6(\mathbb{R}^3) $. Then, considering that $L^6(\mathbb{R}^3) \subset L^6_{\text{loc} }(\mathbb{R}^3)  \subset L^3_{\text{loc} } (\mathbb{R}^3)$, 
thus, by Theorem X.1.1 in \cite{galdi2011introduction} we conclude that $\vu$ and $P$  are in fact a couple of regular functions. 
To continue, let $\psi\in {\mathcal{C}}^\infty_0 (\mathbb{R}^3) $ be a cut-off function such that $0<\psi<1$, $\psi(x)=1$ if $|x|<\frac{1}{2}$, $\psi(x)=0$ if $|x|>1$. Given $R>1$, {\color{black}we define the function $\psi_R$ by $\psi_R(x)=\psi(\frac{|x|}{R})$}: thus, $\psi_R(x)=1$ if $|x|< \frac R 2$ and $\psi_R(x)=0$ if $|x|\geq R$.
 Then, by testing the 3d stationary Navier-Stokes equations \eqref{SNS} with $\psi_R \vu$ and using the fact that $\supp(\psi_R \vu)\subseteq B(R):=B(0,R)$ we obtain \footnote{Note that, since the couple $\vu$ and $\pi$ are regular, the terms involved in the equality   are well-defined.}
\begin{equation}\label{eq. 3}
\int_{B(R)}-\Delta \vu \cdot\left(\psi_R \vu\right)+(\vu \cdot \vn) \vu \cdot\left(\psi_R \vu\right)+\vn \pi \cdot\left(\psi_R \vu\right) d x=0.
\end{equation}
Then, by using the divergence-free condition $\nabla \cdot \vu=0$ and integration by parts, we obtain 
the identity 
\begin{equation*}
\int_{B(R)} \psi_R|\vn \otimes \vu|^2 dx=\int_{B(R)} \Delta \psi_R \frac{|\vu|^2}{2} dx+
\int_{B(R)} \vn \psi_R \cdot\left[\left(\pi+\frac{|\vu|^2}{2}\right) \vu\right]dx.
\end{equation*}
Then, 
since  $\psi_R(x)= 1$ if $|x|< \frac R 2$, we can write
\begin{equation}\label{ineqmother}
\int_{B(R/2)}|\vn \otimes \vu|^2 dx 
 \leq  
\int_{B(R)} \Delta \psi_R \frac{|\vu|^2}{2} dx+ \int_{B(R)} \vn \psi_R \cdot\left[\left(\pi+\frac{|\vu|^2}{2}\right) \vu\right]dx 
=: 
I_1(R)+I_2(R)
. 
\end{equation}
In  the following we will prove 
\begin{equation}\label{LimitsAlphaBeta}
\displaystyle \lim_{R\to +\infty}  |I_1(R)|=\lim_{R\to +\infty}  |I_2(R)|=0.
\end{equation}
{\bf Limit for $I_1(R)$}. For studying the term $I_1(R)$, since the support of $\Delta\psi_R$ is contained in $C(R)$,  the H\"older inequality with 
\begin{equation}\label{variable_exponents_1}
1=\frac{2}{6}+\frac{2}{3}
\end{equation} yields  
\begin{equation}\label{202618aveq32_T1}
|I_1(R)|\leq C
\|\Delta\psi_R\|_{L^{\frac 3 2}(C( R) )}
\|\vu\|_{L^{6}( C( R)  )}^2.
\end{equation}
 In order to control the quantity $\|\Delta \psi_R \|_{L^{  \frac 3 2 }(\mathcal{C}( R) )}$ above,
we write
\begin{equation}\label{EstimationNormeLinfiniTheta_R1}
\|\Delta \psi_R \|_{L^{ \frac 3 2 }(C( R) )}
\leq 
C\|\Delta \psi_R  \|_{L^{\infty}  (C (R) )
}\|1\|_{L^{\frac 3 2 }(C( R) )}.
\end{equation}
Now, considering  the definition of  $\psi_R$,  we get 
$$
\|\Delta \psi_R  \|_{L^{\infty}
( C(R) )
}\leq CR^{-2}$$
 and we obtain
$$\|\Delta \psi_R \|_{L^{ \frac 3 2 }
(\mathcal{C}(  R) )}\leq 
C R^{-2}\|1\|_{L^{ \frac 3 2 }(C(  R) )}
.$$
Now, by stressing the fact   that 
$|
\mathcal{C}( R) 
|=CR^3$ and $R>1$, we obtain
\begin{align}\label{EstimationC1}
\|\Delta \psi_R \|_{L^{\frac 3 2 }(C( R) )}&\leq
  C R^{-2+\frac 3{ \frac 3 2 }} =C
  ,
\end{align}
and thus, we get  
\begin{equation}\label{estimateI12}
|I_1(R)|\leq 
  C  
\|\vu\|_{L^{ 6 }( \mathcal{C}( R)  )}^2.
\end{equation}
Then, since $\|\vu\|_{L^{6}( \mathcal{C}( R)  )} \to 0$ as $R\to +\infty$, we conclude 
$
I_1(R) \to_{R\to +\infty} 0. 
$ 

%%%%%%%%%%%%%%%%%%%%%%%%%%%%%
%%%%%%%%%%%%%%%%%%%%%%%%%%%%%
 \noindent {\bf Limit for $I_2(R)$}. Note that, by considering the definition of $\psi_R$ we know that $\supp( \vn \psi_R)\subset\mathcal{C}( R) $ {and it is radial}. Thus, we can write 
{
\color{black}
\begin{eqnarray}
\left|I_2(R)\right|&=&\left|\int_{B(R)} \vn \psi_R \cdot\left[\left(\pi+\frac{|\vu|^2}{2}\right) \vu\right]dx\right|\notag\\
&=& 
\left|\int_{B(R)} 
\left(\pi+\frac{|\vu|^2}{2}\right) 
\vu_\rho
 \partial_\rho\psi_R 
 dx\right|
\\
& \leq &\frac{1}{2}  
\int_{\mathcal{C}( R) }|\vu|^2
| \vu_\rho |
| \partial_\rho\psi_R |
 dx 
+
 \int_{\mathcal{C}( R) } 
 |\pi |
 | \vu_\rho |
| \partial_\rho\psi_R |  
  dx 
=: I_{21}(R)+I_{22}(R)
.%\label{aug_31T1}
\end{eqnarray}
}

With this at hand in the following we aim to prove 
$$
\displaystyle \lim_{R\to +\infty}I_{21}(R)= \lim_{R\to +\infty} I_{22}(R)=0.
$$
To deal with the term $I_{21}(R)$, by the H\"older inequality with 
$ 
1=\frac{2}{6}+\frac{1}{p} +\frac{1}{q},
$ 
where $p^{-1}+q^{-1}=2/3$
we write 
\begin{align}\label{forheuristic1}
I_{21}(R)= \frac{1}{2}  
\int_{\mathcal{C}( R) }|\vu|^2
| \vu_\rho |
| \partial_\rho\psi_R |
 dx 
&\leq   C
\|\vu \|_{L^{ 6 }(\mathcal{C}( R)  )}^2
\|u_\rho \|_{L^{ p }(\mathcal{C}( R) )}
\|\partial_\rho \psi_R \|_{L^{ q }(\mathcal{C}( R) )}
. 
\end{align}
Since,   $ \| \partial_\rho \psi_R  \|_{L^{\infty}}\leq CR^{-1}$ and $R>1$,  we obtain  
\begin{align}
\| \partial_\rho \psi_R \|_{L^{q}(\mathcal{C}( R) )}
&
\leq 
C\|\partial_\rho \psi_R  \|_{L^{\infty}  }
\|1\|_{L^{ q }(C( R) )}
=
 C R^{-1+\frac 3{q}}  
 ,
\end{align}
Now, considering that $q^{-1}=2/3 - p^{-1}$ and $3/2<p\leq 3$, we have $3 \leq q <+\infty$, and thus we get the uniform bound,  
\begin{equation}
\| \partial_\rho \psi_R \|_{L^{q}(\mathcal{C}( R) )} \leq  C, \quad 
\text{ for all }R> 1.
\end{equation}
By gathering the estimates together, we can write
\begin{equation}%\label{estimateI212}
I_{21}(R)\leq 
  C  
  \|u_\rho \|_{L^{ p }(\mathcal{C}( R) )}
\|\vu \|_{L^{6}( \mathcal{C}( R)  )}^2
\leq 
  C  
  \|u_\rho \|_{L^{ p }(\R^3 )}
\|\vu \|_{L^{6}( \mathcal{C}( R)  )}^2
. 
\end{equation}
Thus, with this information at hand and since $\|u\|_{L^{6}( \mathcal{C}( R)  )} \to 0$ as $R$ evolves to $+\infty$, we conclude $I_{21}(R) \to 0$ as ${R\to +\infty}$.
%%%%%%
Now we analyze the term $I_{22}(R)$. Considering H\"older inequalities with $ \frac{2}{6}+\frac{1}{p } + \frac 1 q =1$, where $p^{-1}+q^{-1}=2/3$,  by arguing in the same manner than before, we  get the estimates 
\begin{align}\label{forheuristic2}
I_{22}(R)
%&
= 
\int_{\mathcal{C}( R) } 
 |\pi |
 | \vu_\rho |
| \partial_\rho\psi_R |  
  dx 
%\\
&\leq 
C 
\|\pi \|_{L^{\frac{6}{2}}(C( R) )}
\|  \partial_\rho\psi_R \|_{L^{q }(C( R) )} 
\| \vu_\rho \|_{L^{p}( C ( R) )} 
\\
&\leq 
 C R^{-1+\frac 3{q }}
\|\pi \|_{L^{\frac{6 }{2}}( C( R) )}\| \vu_\rho \|_{L^{p}( C( R) )}
\\
&=
 C 
\|\pi \|_{L^{\frac{6 }{2}}( \R^3  )}\| \vu_\rho \|_{L^{p}(  C( R) )}
.%\label{estimateI22}
\end{align}
%REVISAR ESTO Y SEGUIR
%
Now,  by using the divergence-free property of $\vu$, we get the following identity for the pressure term
$ \displaystyle{\pi=\sum_{i,j=1}^3\mathcal{R}_i\mathcal{R}_i(u_iu_j)}
,
$ 
where $\mathcal{R}_i$ stands for the usual Riesz transforms.
Then, gathering this relationship with the hypothesis $\vu\in L^{6}(\R^3 )$ and the fact that 
 the Riesz transform are bounded in $L^{r}$ spaces for $r\in (1,\infty)$,  we  conclude   
  \begin{equation}
  \|\pi \|_{L^{ 3}( \R^3 )} \leq 
  C 
    \|u \|^2_{L^{ 6}( \R^3 )  } <+\infty
 ,
  \end{equation}
and then, we can write  
    \begin{align}
I_{22}(R)\leq 
 C  
 \|\vu \|^2_{L^{6}( \R^3 )}
\| \vu_\rho \|_{L^{p}( \R^3 )}
.
\end{align}
Then, considering that   
$\|\vu_\rho\|_{L^{p}( \mathcal{C}( R)   )}\to 0$ as  $R\to +\infty$, 
we get  $ \lim_{R\to +\infty}
I_{22}(R) 
=0$.  Hence, by gathering the limits for $I_{21}(R)$ and $I_{22}(R)$, we deduce the desired limit for $I_{2}(R) $.\\


Thus, by mixing
  the limits for $I_{1}(R)$ and $I_{2}(R)$ with the estimate \eqref{ineqmother}, we conclude
\begin{equation}\label{LimiteSobolev}
\lim_{R\to +\infty}\int_{B(R/2)}|\vn \otimes \vu|^2 dx=\|\vu\|_{\dot{H}^1}=0.
\end{equation}
Then, by considering Sobolev embeddings, we get $\|\vu\|_{L^6}=0$, and in consequence  $\vu= 0$. 
\end{proof}
 

%The proof of this proposition follows exactly the same arguments as Theorem 1.2 in \cite{V2026}, and is therefore omitted here for the sake of brevity.


As noted earlier, the proof of Theorem \ref{thm1.2} follows from a more general Liouville-type result that exploits the flexibility offered by the functional setting of variable-exponent Lebesgue spaces. This result is stated as follows.

%%%%%%%%
\begin{theorem}\label{prop.1}
Let $R_0 > 1$ be fixed, let $u \in \dot{H}^1(\mathbb{R}^3)$ be a solution of \eqref{SNS} and $u_\rho(x)$ be the radial component of such solution in spherical coordinates. Let $p:\mathbb{R}^3 \to [3,6]$ be a continuous and radially decreasing variable exponent such that:
\begin{enumerate}
\item $p(x) = 6$ for all $x \in B(0,R_0)$,
\item there exists a constant $C \in [0, R_0^2)$ such that
$
|p(x) - 3| \le \frac{C}{|x|}, \  \text{for all } |x| \ge R_0^2.
$
\end{enumerate}
If, in addition  $u_\rho \in L^{p(\cdot)}(\mathbb{R}^3)$, then $u \equiv 0$.
\end{theorem}

\begin{proof}%[Proof of Theorem \ref{prop.1}]
 Let  $\vu\in \dot H^1(\R^3) %L^2_{loc}(\R^3 )
   $ be a solution of the 3d stationary Navier-Stokes equations. By considering  Theorem 2.51 in \cite{CRUZ}, we know the inclusions
$$L^{p(\cdot)}(\R^3 )\subset L^{p^-}(\R^3 )+L^{p^+}(\R^3 )
\subset L^{p^-}_{loc}(\R^3 )+L^{p^+}_{loc}(\R^3 )
.$$
Now, by the hypothesis assumed on the variable exponent $p(\cdot)$ we  have $3<9/2\leq p^-\leq p^+\leq 6 $, and  then,  we can deduce  
$$u\in L^{p(\cdot)}(\R^3 ) \subset L^3_{loc}(\R^3 ).$$
Thus, by Theorem X.1.1 in \cite{galdi2011introduction} we conclude that $(\vu,P)$ is in fact a couple of regular functions. By testing the stationary Navier-Stokes equations \eqref{SNS} with $\phi_R \vu$, where $\phi_R $ is the cut-off function defined in the proof of Theorem 1, we get (after some integration by parts)
\begin{equation}\label{ineqmother2}
\int_{B(R/2)}|\vn \otimes \vu|^2 dx 
 \leq  
\int_{B(R)} \Delta \psi_R \frac{|\vu|^2}{2} dx+ \int_{B(R)} \vn \psi_R \cdot\left[\left(\pi+\frac{|\vu|^2}{2}\right) \vu\right]dx 
=: 
J_1(R)+J_2(R)
. 
\end{equation}
In  the following we will prove 
$ \lim_{R\to +\infty}  |J_1(R)|=\lim_{R\to +\infty}  |J_2(R)|=0$.

\noindent {\bf Limit for $J_1(R)$}.
 For studying the term $J_1(R)$ in \eqref{ineqmother2}, the H\"older inequality with 
$%\label{variable_exponents_1_92_exterior_bola}
1=\frac{2}{p(\cdot)}+\frac{1}{q(\cdot)}
$ and Lemma \ref{lemmaLinftyLpvariable}
 yield the  estimates
\footnote{Considering the definition of the cut-off function  $\phi_R$   and Lemma  \ref{lemmaLinftyLpvariable} is straightforward to see that such functions and its partial derivatives belongs to the variable Lebesgue spaces considered here.} 
\begin{align}\label{18aveq32_T1_92_exterior_bola}
|I_1(R)|
&\leq
C\|\Delta \phi_R \|_{L^{q(\cdot)}( \mathcal{C}( R)  )}
\|\vu\|_{L^{p(\cdot)}( \mathcal{C}( R)  )}^2
\\
&\leq 
C
\|\Delta \phi_R  \|_{L^{\infty}  (\mathcal{C}( R)  )}\|1\|_{L^{q_{\mathcal{C}( R)  }(\cdot)}(\mathcal{C}( R)  )}
\|\vu\|_{L^{p(\cdot)}( \mathcal{C}( R)  )}^2
.
\end{align}
Now, considering  the definition of  $\phi_R$,  we get 
$
\|\Delta \phi_R  \|_{L^{\infty}
( \mathcal{C}( R)   )
}\leq CR^{-2}$.
Gathering this  with Lemma \ref{PropositionLpplusminus}
and the fact that
$|
\mathcal{C}( R)  
|=CR^3$  ($R>1$), we obtain the estimates 

\begin{align}\label{EstimationC01_92_exterior_bola}
\|\Delta \phi_R \|_{L^{q_{\mathcal{C}( R)  }(\cdot)}(\mathcal{C}( R)  )}&\leq
 C R^{-2}\max \{|\mathcal{C}( R)  |^{\frac 3 {q_{\mathcal{C}( R)  }^- }}, |\mathcal{C}( R)  |^{\frac 3 {q_{\mathcal{C}( R)  }^+}}\}
 \\ & \leq
  C\max \{R^{-2+\frac 3{q_{\mathcal{C}( R)  }^-}}, R^{-2+\frac 3 {q_{\mathcal{C}( R)  }^+}}\}
\\ & =
  C R^{-2+\frac 3{q_{\mathcal{C}( R)  }^-}} 
 .
\end{align}
Then, we can write
\begin{equation}\label{estimateI1_92_exterior_bola}
|J_1(R)|\leq 
  C R^{-2+\frac 3{q_{\mathcal{C}( R)  }^-}} 
\|\vu\|_{L^{p(\cdot)}( \mathcal{C}( R)  )}^2.
\end{equation}
Then,  considering that $q(\cdot)_{\mathcal{C}( R)  }$ is the conjugate exponent of $p(\cdot)_{\mathcal{C}( R)  }$  and the fact that $\|\vu\|_{L^{p(\cdot)}( \mathcal{C}( R)   )}\to 0$ as  $R\to +\infty$, we conclude 
$$
\lim_{R\to +\infty} J_1(R)  =  0. 
$$

\noindent {\bf Limit for $J_2(R)$}. 
 Note that, by means of  the definition of $\psi_R$ we know that $\supp( \vn \psi_R)\subset\mathcal{C}( R) $ {\color{black}and it is radial}. Thus, we can write 
{
\color{black}
\begin{eqnarray}
\left|I_2(R)\right|
& \leq &\frac{1}{2}  
\int_{\mathcal{C}( R) }|\vu|^2
| \vu_\rho |
| \partial_\rho\psi_R |
 dx 
+
 \int_{\mathcal{C}( R) } 
 |\pi |
 | \vu_\rho |
| \partial_\rho\psi_R |  
  dx 
=: J_{21}(R)+J_{22}(R)
.
\end{eqnarray}
}
With this at hand in the following we aim to prove 
$$
\displaystyle \lim_{R\to +\infty}J_{21}(R)= \lim_{R\to +\infty} J_{22}(R)=0.
$$
To deal with the term $J_{21}(R)$, by the H\"older inequality with 
$ 
1=\frac{2}{6}+\frac{1}{p(\cdot)} +\frac{1}{q(\cdot)},
$ 
where $p^{-1}(\cdot)+q^{-1}(\cdot)=2/3$, 
we write 
\begin{align}
I_{21}(R)= \frac{1}{2}  
\int_{\mathcal{C}( R) }|\vu|^2
| \vu_\rho |
| \partial_\rho\psi_R |
 dx 
&\leq   C
\|\vu \|_{L^{ 6 }(\mathcal{C}( R)  )}^2
\|u_\rho \|_{L^{ p (\cdot) }(\mathcal{C}( R) )}
\|\partial_\rho \psi_R \|_{L^{ q (\cdot) }(\mathcal{C}( R) )}
. 
\end{align}
Since $ \| \partial_\rho \psi_R   \|_{L^{\infty}}\leq CR^{-1}$ and $R>1$, following the same ideas than before, we obtain  
\begin{align}
\| \partial_\rho \psi_R   \|_{L^{q(\cdot)}(\mathcal{C}( R)  )}
&\leq
 C\max \{R^{-1+\frac 3{q_{\mathcal{C}( R)  }^-}}, R^{-1+\frac 3 {q_{\mathcal{C}( R)  }^+}}\} 
=
 C R^{-1+\frac 3{q_{\mathcal{C}( R)  }^-}}
 ,
  \label{EstimationBeta1C_92_exterior_bola}
\end{align}
and then, considering the fact that $q(\cdot)$ is the conjugate exponent of $p(\cdot)$ we get 
\begin{align}\label{estimateI21_92_exterior_bola}
J_{21}(R)
&\leq 
  C R^{-1+\frac 3{q_{\mathcal{C}( R)  }^-}}
\|\vu \|_{L^{ 6 }(\mathcal{C}( R)  )}^2
\|u_\rho \|_{L^{ p (\cdot) }(\mathcal{C}( R) )}
\\ &\leq 
  C R^{1-\frac 3{p_{\mathcal{C}( R)  }^+}}
\|\vu \|_{L^{ 6 }(\mathcal{C}( R)  )}^2
\|u_\rho \|_{L^{ p (\cdot) }(\R^3 )}
. 
\end{align}
Then, since  
$\|\vu\|_{L^{6}( \mathcal{C}( R)   )}\to 0$ as  $R\to +\infty$, by considering Lemma \ref{prop:main_limit_thmover_new} we obtain
$$\lim_{R\to +\infty} R^{1-\frac 3{p_{\mathcal{C}( R)  }^+}}
\|\vu \|_{L^{ 6 }(\mathcal{C}( R)  )}^2
\|u_\rho \|_{L^{ p (\cdot) }(\R^3 )}
= 0,$$
and therefore $$
\lim_{R\to +\infty} J_{21}(R) = 0.$$ 
%%%%%%
Now we analyze the term $J_{22}(R)$. Considering H\"older inequalities with $ \frac{2}{6}+\frac{1}{p (\cdot) } + \frac 1 {q(\cdot)} =1$, where $p^{-1}(\cdot)+q^{-1}(\cdot)=2/3$,  by arguing in the same manner than before, we  get the estimates 
\begin{align}
J_{22}(R)
%&
= 
\int_{\mathcal{C}( R) } 
 |\pi |
 | \vu_\rho |
| \partial_\rho\psi_R |  
  dx 
%\\
&\leq 
C 
\|\pi \|_{L^{\frac{6}{2}}(C( R) )}
\|  \partial_\rho\psi_R \|_{L^{q (\cdot)  }(C( R) )} 
\| \vu_\rho \|_{L^{p(\cdot) }( C ( R) )} 
\\
&\leq 
 C  R^{-1-\frac 3{q_{\mathcal{C}( R)  }^+}}
\|\pi \|_{L^{\frac{6 }{2}}( C( R) )}\| \vu_\rho \|_{L^{p(\cdot)}( C( R) )}
\\
&=
 C R^{1-\frac 3{p_{\mathcal{C}( R)  }^+}}
\|\pi \|_{L^{\frac{6 }{2}}( C( R) )}\| \vu_\rho \|_{L^{p(\cdot)}(\mathcal{C}( R)  )}
.%\label{estimateI22}
\end{align}
Then, considering the identity for the pressure term $\pi=\sum_{i,j=1}^3\mathcal{R}_i\mathcal{R}_i(u_iu_j)$  and the fact that the Riesz transforms are bounded in the classical Lebesgue spaces involved here, we can write 
  \begin{equation}
  \|\pi \|_{L^{ 3}( \R^3 )} \leq 
  C 
    \|u \|^2_{L^{ 6}( \R^3 )  } 
  ,
  \end{equation}
and hence  we conclude
  \begin{align}
J_{22}(R)\leq 
 C R^{1-\frac 3{p_{\mathcal{C}( R)  }^+}}
 \|\vu \|^2_{L^{6}( \R^3 )}
\| \vu_\rho \|_{L^{p(\cdot)}(\mathcal{C}( R)  )}
.
\end{align}
Thus, considering that 
$\|\vu_\rho\|_{L^{p(\cdot)}( \mathcal{C}( R)   )}\to 0$ as  $R\to +\infty$, by Lemma \ref{prop:main_limit_thmover_new} we get
 $$
\lim_{R\to +\infty} J_{22}(R) = 0.$$ 
Gathering all the information of these limits  with (\ref{ineqmother2}), we obtain
\begin{equation} 
\lim_{R\to +\infty}\int_{B(R/2)}|\vn \otimes \vu|^2 dx=\|\vu\|_{\dot{H}^1(\R^3)}=0.
\end{equation}
Thus, 
considering Sobolev embeddings we  get $\|\vu\|_{L^6}=0$ and thus  $\vu\equiv 0$. 
%%
\end{proof}

%%%%%%%%%%%%%%%%%%%%%%%
%\newpage
\begin{proof}[Proof of Theorem \ref{thm1.2}]
Consider the variable exponent $p(\cdot)= 3 + \eta(\cdot)$. 
Then, for each $x\in \R^3$, we have $p(x)\geq 3$ and it is in fact a radial function. Furthermore,  given $|x|=R_0$, we can write 
    \[
p(R_0)=     3 +3 \frac{R_0}{R_0} = 3+ 3 = 6.
    \]
This proves that the variable exponent  $p(\cdot)$ is in fact  continuous.
To continue,  given $ |x| > R_0 $,  $ \frac{R_0}{|x|} $ is  a decreasing function, therefore $ p(\cdot) $ is also  decreasing function. Now, given  $ |x| \geq R_0^2 $, we get  
\[
\left| p(x) -  3 \right| =
 \left| 3 \frac{R_0}{|x|} \right| = C |x|^{-1},
\]
with $ C =  3 R_0 <R_0^2$ ($R_0>3$ by assumption). Provided with this, we stress the fact that this variable exponent $p(\cdot)$ fulfill the hypothesis of Theorem \ref{prop.1}, therefore if the radial  component of the velocity fielf $u_\rho$ belongs to $L^{p(\cdot)} (\R^3)$, we have $u\equiv 0$.
\end{proof}


%\paragraph{\bf Acknowledgements} The author warmly thanks  Hedong Hou and Alexey Cheskidov  for their valuable advice and insightful comments.
 
%\paragraph{\bf Datasets} Is important to mention that data sharing does not apply to this article as no datasets were generated or analyzed during the current research study.
 
%\paragraph{\bf  Conflict of interest} The author declares no conflict of interest and confirms being the sole contributor to this paper.

%------------------
%------ end of document 
%------------------
%\newpage
\begin{thebibliography}{99}
%\addcontentsline{toc}{section}{References}
%

\bibitem{Chae23}
{\sc D.~Chae}, {\em Anisotropic Liouville type theorem for the stationary Navier–Stokes
equations in $\R^3$}, 
Appl. Math. Lett. 142, Paper No. 108655, 5 pp. (2023).
%
\bibitem{ChaeWolf}
{\sc D.~Chae and J.~Wolf}, {\em On {L}iouville type theorems for the steady
  {N}avier-{S}tokes equations in {$\mathbb{R}^3$}}, J. Differential Equations,
  261 (2016), pp.~5541--5560.
%
\bibitem{CV2023} {\sc D. Chamorro and  G. Vergara-Hermosilla}, {\em  Liouville type theorems for stationary Navier-Stokes equations with Lebesgue spaces of variable exponent}, Documenta Mathematica (2025).
%
\bibitem{CRUZ}
{\sc D.~V. Cruz-Uribe and A.~Fiorenza}, {\em Variable {L}ebesgue spaces:
Foundations and harmonic analysis}, Springer Science \& Business Media, 2013.
%
\bibitem{Diening_Libro}
{\sc L.~Diening, P.~Harjulehto, P.~H{\"a}st{\"o} and M.~Ruzicka}, {\em
Lebesgue and Sobolev spaces with variable exponents}, Springer, 2011.
%
\bibitem{galdi2011introduction}
{\sc G.~Galdi}, {\em An introduction to the mathematical theory of the
{N}avier-{S}tokes equations: Steady-state problems}, Springer Science \&
Business Media, 2011.
%
\bibitem{GW78}
{\sc D. Gilbarg and  H.F.  Weinberger}, {\em Asymptotic properties of steady plane solutions
of the Navier-Stokes equations with bounded Dirichlet integral}, Ann. Scuola
Norm. Super. Pisa. CI. Sci. 5(4)  (1978), 381–404.
%
\bibitem{lemarie2016navier}
{\sc P.~G. Lemari{\'e}-Rieusset},
{\em The {N}avier-{S}tokes problem in the 21st century}, CRC press, 2016.
%
%
\bibitem{Lerner26}
{\sc N.~
Lerner},   {\em Wiener Algebras Methods for Liouville Theorems on the Stationary Navier-Stokes System}, arXiv preprint arXiv:2601.13916  (2026).
%
\bibitem{Ser2016}  
\textsc{G. Seregin}, \emph{A Liouville type theorem for steady-state Navier-Stokes equations}.  
J. {\'E}.D.P., Expos{\'e} no IX, (2016).
%
\bibitem{V2026}
 {\sc G. Vergara-Hermosilla}, \emph{Liouville theorems above the critical $9/2$ Threshold for stationary Navier-Stokes Equations}, arXiv preprint arXiv:2604.06527, 2026.
%
\bibitem{ZZ25}
 {\sc  Z. Zhang and  Q. Zu}, 
  {\em Two improved anisotropic Liouville type theorems for the stationary 3D Navier–Stokes equations}, 
Arch. Math. 124 (2025), 571–582.

\end{thebibliography}

 \end{document}
